\MathBookSourcePage{322}
\subsection{弱形式与 Korn 不等式}
\label{sec:ch11-weak-form-korn}
\setcounter{thmcount}{0}
\setcounter{equation}{0}

假设 $\chvec u\in\chvec H^2(\Omega)$ 满足~\eqref{eq:ch11-static-elasticity-equation}. 则对任意 $\chvec v\in\chvec H^1(\Omega)$ 且 $\chvec v|_{\Gamma_1}=0$, 分部积分(假设可以进行)得到
\begin{align*}
\MathBookFormulaID{formula-ch11-elasticity-integration-by-parts}
\int_\Omega\chvec f\cdot\chvec v\,dx
&=-\int_\Omega\chvec{\operatorname{div}}\chten\sigma(\chvec u)\cdot\chvec v\,dx\\
&=\int_\Omega\{2\mu\chten\epsilon(\chvec u)+
\lambda\operatorname{tr}(\chten\epsilon(\chvec u))\chten\delta\}:
\chten{\operatorname{grad}}\chvec v\,dx
-\int_{\Gamma_2}\chten\sigma(\chvec u)\chvec\nu\cdot\chvec v\,ds.
\end{align*}
因此, 由~\eqref{eq:ch11-static-elasticity-equation}--\eqref{eq:ch11-traction-boundary-condition} 得到如下弱形式:
求 $\chvec u\in\chvec H^1(\Omega)$, 使得 $\chvec u|_{\Gamma_1}=\chvec g$, 并且
\MBSourceItem{1}
\begin{equation}
\MathBookFormulaID{formula-ch11-elasticity-weak-problem}
\label{eq:ch11-elasticity-weak-problem}
a(\chvec u,\chvec v)=\int_\Omega\chvec f\cdot\chvec v\,dx
+\int_{\Gamma_2}\chvec t\cdot\chvec v\,ds
\qquad\forall\chvec v\in\chvec V,
\end{equation}
其中
\MBSourceItem{2}
\begin{equation}
\MathBookFormulaID{formula-ch11-elasticity-bilinear-form}
\label{eq:ch11-elasticity-bilinear-form}
\begin{split}
a(\chvec u,\chvec v)
&:=\int_\Omega\left\{2\mu\chten\epsilon(\chvec u):\chten\epsilon(\chvec v)
+\lambda\operatorname{div}\chvec u\operatorname{div}\chvec v\right\}\,dx,\\
\chvec V&:=\{\chvec v\in\chvec H^1(\Omega):\chvec v|_{\Gamma_1}=0\}.
\end{split}
\end{equation}

首先必须问: 问题~\eqref{eq:ch11-elasticity-weak-problem} 是否有唯一解? 换言之, 需要建立有界双线性形式 $a(\cdot,\cdot)$ 在 $\chvec V$ 上的强制性. 我们从一个关于欠定边值问题的引理开始(注意, 其解不唯一), 本章及第~\MBUnresolvedReference{chap:ch12-mixed-methods}{12} 章将频繁使用它. 其证明使用了超出第~\MBUnresolvedReference{chap:ch01}{1} 章所发展的 Sobolev 空间技术, 但为完整起见给出证明梗概. 本节其余部分假设 $\Omega$ 或为多边形, 或 $\partial\Omega$ 光滑.

\MathBookSourcePage{323}
\MBSourceItem{3}
\begin{Lemma}
\label{lem:ch11-divergence-right-inverse}
存在正常数 $C$, 使得对所有 $p\in L^2(\Omega)$, 都存在 $\chvec v\in\chvec H^1(\Omega)$ 满足
\[
\MathBookFormulaID{formula-ch11-divergence-right-inverse-equation}
\operatorname{div}\chvec v=p
\]
以及
\[
\MathBookFormulaID{formula-ch11-divergence-right-inverse-bound}
\|\chvec v\|_{\chvec H^1(\Omega)}\leq C\|p\|_{L^2(\Omega)}.
\]
此外, 若 $p$ 满足 $\int_\Omega p\,dx=0$, 则可取 $\chvec v\in\mathring{\chvec H}^1(\Omega)$.
\end{Lemma}

\begin{Proof}
对于边界光滑的 $\Omega$ 和多边形 $\Omega$, 该引理均成立. 这里给出光滑边界情形的证明, 多边形情形参见 Girault 与 Raviart (1986), Arnold, Scott 与 Vogelius (1988). 存在唯一的 $w\in H^2(\Omega)$ 满足
\[
\MathBookFormulaID{formula-ch11-dirichlet-poisson-for-divergence}
-\Delta w=p\quad\text{in }\Omega,
\qquad w=0\quad\text{on }\partial\Omega.
\]
由椭圆正则性(参见第~\MBUnresolvedReference{sec:ch05-biharmonic}{5.5} 节),
\[
\MathBookFormulaID{formula-ch11-dirichlet-poisson-regularity}
\|w\|_{H^2(\Omega)}\leq C_2\|p\|_{L^2(\Omega)}.
\]
于是 $\chvec v=-\chvec{\operatorname{grad}}w$ 满足引理条件.

现在假设 $\int_\Omega p\,dx=0$. 则存在唯一的 $w\in H^2(\Omega)$ 满足
\MBSourceItem{4}
\begin{equation}
\MathBookFormulaID{formula-ch11-neumann-poisson-for-divergence}
\label{eq:ch11-neumann-poisson-for-divergence}
-\Delta w=p\quad\text{in }\Omega,
\qquad
\frac{\partial w}{\partial\chvec\nu}=0\quad\text{on }\partial\Omega,
\qquad
\int_\Omega w\,dx=0
\end{equation}
(参见第~\MBUnresolvedReference{sec:ch05-neumann}{5.2} 节). 由椭圆正则性,
\[
\MathBookFormulaID{formula-ch11-neumann-poisson-regularity}
\|w\|_{H^2(\Omega)}\leq C_2\|p\|_{L^2(\Omega)}.
\]
\MathBookSourcePage{324}
令 $\chvec v_1=-\chvec{\operatorname{grad}}w$. 则 $\chvec v_1\in\chvec H^1(\Omega)$,
\MBSourceItem{5}
\begin{equation}
\MathBookFormulaID{formula-ch11-v1-divergence}
\label{eq:ch11-v1-divergence}
\operatorname{div}\chvec v_1=p
\end{equation}
且
\MBSourceItem{6}
\begin{equation}
\MathBookFormulaID{formula-ch11-v1-bound}
\label{eq:ch11-v1-bound}
\|\chvec v_1\|_{\chvec H^1(\Omega)}\leq C_2\|p\|_{L^2(\Omega)}.
\end{equation}
记 $\chvec\nu$ 为 $\partial\Omega$ 上的单位外法向量. 由~\eqref{eq:ch11-neumann-poisson-for-divergence} 可知
\MBSourceItem{7}
\begin{equation}
\MathBookFormulaID{formula-ch11-v1-normal-zero}
\label{eq:ch11-v1-normal-zero}
\chvec v_1|_{\partial\Omega}\cdot\chvec\nu
=-\chvec{\operatorname{grad}}w|_{\partial\Omega}\cdot\chvec\nu=0.
\end{equation}
令 $\chvec\tau$ 为正向单位切向量. 迹定理(参见 Adams 1975)蕴含存在 $\psi\in H^2(\Omega)$, 使得
\[
\MathBookFormulaID{formula-ch11-psi-boundary-data}
\psi|_{\partial\Omega}=0,
\qquad
\frac{\partial\psi}{\partial\chvec\nu}\bigg|_{\partial\Omega}
=\chvec v_1|_{\partial\Omega}\cdot\chvec\tau,
\]
且对某个正常数 $C_\tau$,
\MBSourceItem{8}
\begin{equation}
\MathBookFormulaID{formula-ch11-psi-bound}
\label{eq:ch11-psi-bound}
\|\psi\|_{H^2(\Omega)}\leq C_\tau\|\chvec v_1\|_{\chvec H^1(\Omega)}.
\end{equation}
令 $\chvec v_2=\chvec{\operatorname{curl}}\psi$. 则
\MBSourceItem{9}
\begin{equation}
\MathBookFormulaID{formula-ch11-v2-normal-zero}
\label{eq:ch11-v2-normal-zero}
\chvec v_2|_{\partial\Omega}\cdot\chvec\nu
=\chvec{\operatorname{grad}}\psi|_{\partial\Omega}\cdot\chvec\tau=0,
\end{equation}
且
\MBSourceItem{10}
\begin{equation}
\MathBookFormulaID{formula-ch11-v2-tangential-cancellation}
\label{eq:ch11-v2-tangential-cancellation}
\chvec v_2|_{\partial\Omega}\cdot\chvec\tau
=-\chvec{\operatorname{grad}}\psi|_{\partial\Omega}\cdot\chvec\nu
=-\chvec v_1|_{\partial\Omega}\cdot\chvec\tau.
\end{equation}
结合~\eqref{eq:ch11-v1-normal-zero}, \eqref{eq:ch11-v2-normal-zero} 和~\eqref{eq:ch11-v2-tangential-cancellation}, 得到
\MBSourceItem{11}
\begin{equation}
\MathBookFormulaID{formula-ch11-boundary-cancellation}
\label{eq:ch11-boundary-cancellation}
\chvec v_2|_{\partial\Omega}=-\chvec v_1|_{\partial\Omega}.
\end{equation}
令 $\chvec v=\chvec v_1+\chvec v_2$. 由~\eqref{eq:ch11-v1-divergence} 有
$\operatorname{div}\chvec v=\operatorname{div}\chvec v_1=p$, 由~\eqref{eq:ch11-boundary-cancellation} 有 $\chvec v|_{\partial\Omega}=0$, 并且
\begin{align*}
\MathBookFormulaID{formula-ch11-divergence-right-inverse-final-bound}
\|\chvec v\|_{\chvec H^1(\Omega)}
&\leq\|\chvec v_1\|_{\chvec H^1(\Omega)}+\|\chvec v_2\|_{\chvec H^1(\Omega)}\\
&\leq\|\chvec v_1\|_{\chvec H^1(\Omega)}+C\|\psi\|_{H^2(\Omega)}
\leq C\|\chvec v_1\|_{\chvec H^1(\Omega)}
\leq C\|p\|_{L^2(\Omega)}.
\end{align*}
\end{Proof}

令 $\widehat{\chvec H}^k(\Omega)$ 定义为
\[
\MathBookFormulaID{formula-ch11-hat-vector-sobolev-space}
\widehat{\chvec H}^k(\Omega):=
\left\{\chvec v\in\chvec H^k(\Omega):
\int_\Omega\chvec v\,dx=\chvec0,
\ \int_\Omega\operatorname{rot}\chvec v\,dx=0\right\}.
\]
当 $k\geq1$ 时, 它们是 $\chvec H^k(\Omega)$ 的闭子空间.

\MathBookSourcePage{325}
\MBSourceItem{12}
\begin{Theorem}[第二 Korn 不等式]
\label{thm:ch11-second-korn}
存在正常数 $C$, 使得
\MBSourceItem{13}
\begin{equation}
\MathBookFormulaID{formula-ch11-second-korn-inequality}
\label{eq:ch11-second-korn-inequality}
\|\chten\epsilon(\chvec v)\|_{\chten L^2(\Omega)}
\geq C\|\chvec v\|_{\chvec H^1(\Omega)}
\qquad\forall\chvec v\in\widehat{\chvec H}^1(\Omega).
\end{equation}
\end{Theorem}

\begin{Proof}
令 $\chvec v\in\widehat{\chvec H}^1(\Omega)$. 由于 $\int_\Omega\operatorname{rot}\chvec v\,dx=0$, 引理~\ref{lem:ch11-divergence-right-inverse} 蕴含存在 $\chvec w\in\mathring{\chvec H}^1(\Omega)$, 使得
\MBSourceItem{14}
\begin{equation}
\MathBookFormulaID{formula-ch11-w-divergence-rotation}
\label{eq:ch11-w-divergence-rotation}
\operatorname{div}\chvec w=\operatorname{rot}\chvec v,
\qquad
\|\chvec w\|_{\chvec H^1(\Omega)}\leq C_1\|\operatorname{rot}\chvec v\|_{L^2(\Omega)}.
\end{equation}
于是, 由~\eqref{eq:ch11-symmetric-gradient-identity},
\begin{align*}
\MathBookFormulaID{formula-ch11-second-korn-proof-identity}
\int_\Omega\chten\epsilon(\chvec v):
(\chten{\operatorname{grad}}\chvec v-\chten{\operatorname{curl}}\chvec w)\,dx
&=\int_\Omega\left(\chten{\operatorname{grad}}\chvec v-
\frac12(\operatorname{rot}\chvec v)\chten\chi\right):
(\chten{\operatorname{grad}}\chvec v-\chten{\operatorname{curl}}\chvec w)\,dx\\
&=\|\chten{\operatorname{grad}}\chvec v\|_{\chten L^2(\Omega)}^2
-\int_\Omega\chten{\operatorname{grad}}\chvec v:
\chten{\operatorname{curl}}\chvec w\,dx\\
&\quad-\frac12\int_\Omega(\operatorname{rot}\chvec v)
(\chten\chi:\chten{\operatorname{grad}}\chvec v-
\chten\chi:\chten{\operatorname{curl}}\chvec w)\,dx\\
&=\|\chten{\operatorname{grad}}\chvec v\|_{\chten L^2(\Omega)}^2
-\frac12\int_\Omega(\operatorname{rot}\chvec v)
(\operatorname{rot}\chvec v-\operatorname{div}\chvec w)\,dx\\
&=\|\chten{\operatorname{grad}}\chvec v\|_{\chten L^2(\Omega)}^2.
\end{align*}
因此, Schwarz 不等式和~\eqref{eq:ch11-w-divergence-rotation} 蕴含
\MBSourceItem{15}
\begin{equation}
\MathBookFormulaID{formula-ch11-second-korn-gradient-bound}
\label{eq:ch11-second-korn-gradient-bound}
\|\chten{\operatorname{grad}}\chvec v\|_{\chten L^2(\Omega)}^2
\leq\|\chten\epsilon(\chvec v)\|_{\chten L^2(\Omega)}
\|\chten{\operatorname{grad}}\chvec v-
\chten{\operatorname{curl}}\chvec w\|_{\chten L^2(\Omega)}
\leq C\|\chten\epsilon(\chvec v)\|_{\chten L^2(\Omega)}
\|\chvec v\|_{\chvec H^1(\Omega)}.
\end{equation}
定理由~\eqref{eq:ch11-second-korn-gradient-bound} 和 Friedrichs 不等式~\MBUnresolvedReference{eq:ch04-source-4-3-15}{4.3.15} 得出, 因为 $\int_\Omega\chvec v\,dx=0$.
\end{Proof}

\MBSourceItem{16}
\begin{Theorem}[Korn 不等式]
\label{thm:ch11-korn}
存在正常数 $\alpha$, 使得
\MBSourceItem{17}
\begin{equation}
\MathBookFormulaID{formula-ch11-korn-inequality}
\label{eq:ch11-korn-inequality}
\|\chten\epsilon(\chvec v)\|_{\chten L^2(\Omega)}^2
+\|\chvec v\|_{\chvec L^2(\Omega)}^2
\geq\alpha\|\chvec v\|_{\chvec H^1(\Omega)}^2
\qquad\forall\chvec v\in\chvec H^1(\Omega).
\end{equation}
\end{Theorem}

\begin{Proof}
首先注意 $\chvec H^1(\Omega)=\widehat{\chvec H}^1(\Omega)\oplus\chvec{\mathrm{RM}}$(参见习题~\ref{exr:ch11-h1-rm-decomposition}), 其中``$\oplus$''只按代数意义理解(它们在 $\chvec H^1(\Omega)$ 中不正交). 因此, 给定任意 $\chvec v\in\chvec H^1(\Omega)$, 存在唯一的一对
$(\chvec z,\chvec w)\in\widehat{\chvec H}^1(\Omega)\times\chvec{\mathrm{RM}}$, 使得
$\chvec v=\chvec z+\chvec w$.

\MathBookSourcePage{326}
特别地, $\chvec w$ 具有~\eqref{eq:ch11-rigid-motions-space} 中的形式, 其中
\MBSourceItem{18}
\begin{equation}
\MathBookFormulaID{formula-ch11-rigid-motion-coefficients}
\label{eq:ch11-rigid-motion-coefficients}
b:=-\frac1{2\operatorname{meas}(\Omega)}
\int_\Omega\operatorname{rot}\chvec v\,dx,
\qquad
\chvec c:=\frac1{\operatorname{meas}(\Omega)}
\int_\Omega\{\chvec v(x)-b(x_2,-x_1)^t\}\,dx.
\end{equation}
因此, $\|\chvec w\|_{\chvec H^1(\Omega)}\leq C\|\chvec v\|_{\chvec H^1(\Omega)}$, 其中 $C$ 只依赖于 $\operatorname{meas}(\Omega)$. 由三角不等式, 存在正常数 $C_1$(参见习题~\ref{exr:ch11-closed-subspace-angle}), 使得
\MBSourceItem{19}
\begin{equation}
\MathBookFormulaID{formula-ch11-decomposition-norm-lower-bound}
\label{eq:ch11-decomposition-norm-lower-bound}
C_1\left(\|\chvec z\|_{\chvec H^1(\Omega)}+
\|\chvec w\|_{\chvec H^1(\Omega)}\right)
\leq\|\chvec v\|_{\chvec H^1(\Omega)}.
\end{equation}

用反证法证明定理. 若~\eqref{eq:ch11-korn-inequality} 对任何正常数 $C$ 都不成立, 则存在序列 $\{\chvec v_n\}\subseteq\chvec H^1(\Omega)$, 使得
\MBSourceItem{20}
\begin{equation}
\MathBookFormulaID{formula-ch11-korn-contradiction-normalization}
\label{eq:ch11-korn-contradiction-normalization}
\|\chvec v_n\|_{\chvec H^1(\Omega)}=1
\end{equation}
且
\MBSourceItem{21}
\begin{equation}
\MathBookFormulaID{formula-ch11-korn-contradiction-smallness}
\label{eq:ch11-korn-contradiction-smallness}
\|\chten\epsilon(\chvec v_n)\|_{\chten L^2(\Omega)}^2
+\|\chvec v_n\|_{\chvec L^2(\Omega)}^2<\frac1n.
\end{equation}
对每个 $n$, 令 $\chvec v_n=\chvec z_n+\chvec w_n$, 其中
$\chvec z_n\in\widehat{\chvec H}^1(\Omega)$ 且 $\chvec w_n\in\chvec{\mathrm{RM}}$. 于是
\[
\MathBookFormulaID{formula-ch11-zn-strain-smallness}
\|\chten\epsilon(\chvec z_n)\|_{\chten L^2(\Omega)}
=\|\chten\epsilon(\chvec v_n)\|_{\chten L^2(\Omega)}<\frac1n.
\]
第二 Korn 不等式蕴含 $\chvec z_n\to0$ 于 $\chvec H^1(\Omega)$. 由~\eqref{eq:ch11-decomposition-norm-lower-bound} 和~\eqref{eq:ch11-korn-contradiction-normalization}, $\{\chvec w_n\}$ 是 $\chvec H^1(\Omega)$ 中的有界序列. 由于 $\chvec{\mathrm{RM}}$ 是三维的, $\{\chvec w_n\}$ 有收敛子序列 $\{\chvec w_{n_j}\}$. 于是子序列 $\{\chvec v_{n_j}=\chvec z_{n_j}+\chvec w_{n_j}\}$ 在 $\chvec H^1(\Omega)$ 中收敛到某个 $\chvec v\in\chvec{\mathrm{RM}}$. 由~\eqref{eq:ch11-korn-contradiction-normalization} 和~\eqref{eq:ch11-korn-contradiction-smallness}, 得到
\[
\MathBookFormulaID{formula-ch11-korn-contradiction-limit}
\|\chvec v\|_{\chvec H^1(\Omega)}=1,
\qquad
\|\chvec v\|_{\chvec L^2(\Omega)}=0,
\]
矛盾.
\end{Proof}

\MBSourceItem{22}
\begin{Corollary}
\label{cor:ch11-korn-subspace}
设 $\chvec V$ 由~\eqref{eq:ch11-elasticity-bilinear-form} 定义且 $\operatorname{meas}(\Gamma_1)>0$. 则存在正常数 $C$, 使得
\MBSourceItem{23}
\begin{equation}
\MathBookFormulaID{formula-ch11-korn-subspace-inequality}
\label{eq:ch11-korn-subspace-inequality}
\|\chten\epsilon(\chvec v)\|_{\chten L^2(\Omega)}
\geq C\|\chvec v\|_{\chvec H^1(\Omega)}
\qquad\forall\chvec v\in\chvec V.
\end{equation}
\end{Corollary}

\MathBookSourcePage{327}
\begin{Proof}
同样的反证法证明存在某个 $\chvec v\in\chvec{\mathrm{RM}}\cap\chvec V$, 使得
\MBSourceItem{24}
\begin{equation}
\MathBookFormulaID{formula-ch11-korn-subspace-contradiction}
\label{eq:ch11-korn-subspace-contradiction}
\|\chvec v\|_{\chvec H^1(\Omega)}=1.
\end{equation}
但 $\chvec{\mathrm{RM}}$ 中唯一满足 $\chvec v|_{\Gamma_1}=0$ 的 $\chvec v$ 是 $\chvec v=0$, 因为 $\operatorname{meas}(\Gamma_1)>0$(参见习题~\ref{exr:ch11-rigid-motion-zero-set}). 这与~\eqref{eq:ch11-korn-subspace-contradiction} 矛盾.
\end{Proof}

在 $\Gamma_2=\varnothing$ 的情形, 立即得到以下结果.

\MBSourceItem{25}
\begin{Corollary}[第一 Korn 不等式]
\label{cor:ch11-first-korn}
存在正常数 $C$, 使得
\MBSourceItem{26}
\begin{equation}
\MathBookFormulaID{formula-ch11-first-korn-inequality}
\label{eq:ch11-first-korn-inequality}
\|\chten\epsilon(\chvec v)\|_{\chten L^2(\Omega)}
\geq C\|\chvec v\|_{\chvec H^1(\Omega)}
\qquad\forall\chvec v\in\mathring{\chvec H}^1(\Omega).
\end{equation}
\end{Corollary}

\MBSourceItem{27}
\begin{Remark}
\label{rem:ch11-korn-generality}
Korn 不等式实际上对具有 Lipschitz 边界的 $\Omega\subseteq\mathbb R^n$ 成立. 一般情形的证明参见 Duvaut 与 Lions (1972), Nitsche (1981). 它们也可推广到分片 $H^1$ 向量场(Brenner 2004).
\end{Remark}

现在注意, 当 $\operatorname{meas}(\Gamma_1)>0$ 时, $a(\cdot,\cdot)$ 在 $\chvec V$ 上的强制性立即由推论~\ref{cor:ch11-korn-subspace} 得出, 从而得到~\eqref{eq:ch11-elasticity-weak-problem} 的唯一可解性.

\MBSourceItem{28}
\begin{Theorem}
\label{thm:ch11-mixed-boundary-unique-solvability}
假设 $\chvec f\in\chvec H^{-1}(\Omega)$, $\chvec g=\chvec w|_{\Gamma_1}$, 其中 $\chvec w\in\chvec H^1(\Omega)$, $\chvec t\in\chvec L^2(\Gamma_2)$, 且 $\operatorname{meas}(\Gamma_1)>0$. 则变分问题~\eqref{eq:ch11-elasticity-weak-problem} 有唯一解.
\end{Theorem}

\begin{Proof}
令 $\chvec u^*=\chvec u-\chvec w$. 则~\eqref{eq:ch11-elasticity-weak-problem} 等价于求 $\chvec u^*\in\chvec V$, 使得对所有 $\chvec v\in\chvec V$,
\MBSourceItem{29}
\begin{equation}
\MathBookFormulaID{formula-ch11-shifted-elasticity-functional}
\label{eq:ch11-shifted-elasticity-functional}
\begin{split}
a(\chvec u^*,\chvec v)
&=\int_\Omega\chvec f\cdot\chvec v\,dx
+\int_{\Gamma_2}\chvec t\cdot\chvec v\,ds\\
&\quad-\int_\Omega\{2\mu\chten\epsilon(\chvec w):\chten\epsilon(\chvec v)
+\lambda\operatorname{div}\chvec w\operatorname{div}\chvec v\}\,dx
=:F(\chvec v),
\end{split}
\end{equation}
其中, 由 $\chvec f$, $\chvec t$ 和 $\chvec w$ 的条件, $F\in\chvec V'$. 因此, 由定理~\MBUnresolvedReference{thm:ch02-source-2-5-6}{2.5.6}, 式~\eqref{eq:ch11-shifted-elasticity-functional} 有唯一解, 这蕴含~\eqref{eq:ch11-elasticity-weak-problem} 也有唯一解.
\end{Proof}

注意在纯牵引问题的情形, $\chvec V$ 变为 $\chvec H^1(\Omega)$. 由于齐次纯牵引问题有非平凡核 $\chvec{\mathrm{RM}}$, 在问题~\eqref{eq:ch11-elasticity-weak-problem} 可解之前, $\chvec f$ 和 $\chvec t$ 必须满足某些相容性条件.

\MathBookSourcePage{328}
\MBSourceItem{30}
\begin{Theorem}
\label{thm:ch11-pure-traction-solvability}
假设 $\chvec f\in\chvec L^2(\Omega)$ 且 $\chvec t\in\chvec L^2(\Gamma)$. 则变分问题
\MBSourceItem{31}
\begin{equation}
\MathBookFormulaID{formula-ch11-pure-traction-variational-problem}
\label{eq:ch11-pure-traction-variational-problem}
\text{求 }\chvec u\in\chvec H^1(\Omega)\text{, 使得}\qquad
a(\chvec u,\chvec v)=\int_\Omega\chvec f\cdot\chvec v\,dx
+\int_\Gamma\chvec t\cdot\chvec v\,ds
\quad\forall\chvec v\in\chvec H^1(\Omega)
\end{equation}
可解, 当且仅当满足如下相容性条件:
\MBSourceItem{32}
\begin{equation}
\MathBookFormulaID{formula-ch11-pure-traction-compatibility}
\label{eq:ch11-pure-traction-compatibility}
\int_\Omega\chvec f\cdot\chvec v\,dx
+\int_\Gamma\chvec t\cdot\chvec v\,ds=0
\qquad\forall\chvec v\in\chvec{\mathrm{RM}}.
\end{equation}
当~\eqref{eq:ch11-pure-traction-variational-problem} 可解时, 在 $\widehat{\chvec H}^1(\Omega)$ 中存在唯一解.
\end{Theorem}

\begin{Proof}
必要性: 若~\eqref{eq:ch11-pure-traction-variational-problem} 可解, 则由~\eqref{eq:ch11-rigid-motion-strain-zero},
\[
\MathBookFormulaID{formula-ch11-pure-traction-necessity}
\int_\Omega\chvec f\cdot\chvec v\,dx+
\int_\Gamma\chvec t\cdot\chvec v\,ds
=a(\chvec u,\chvec v)=0
\qquad\forall\chvec v\in\chvec{\mathrm{RM}}.
\]

充分性: 假设~\eqref{eq:ch11-pure-traction-compatibility} 成立. 由第二 Korn 不等式~\eqref{eq:ch11-second-korn-inequality} 和定理~\MBUnresolvedReference{thm:ch02-source-2-5-6}{2.5.6}, 存在唯一的 $\chvec u^*\in\widehat{\chvec H}^1(\Omega)$, 使得
\[
\MathBookFormulaID{formula-ch11-pure-traction-restricted-solution}
a(\chvec u^*,\chvec v)=\int_\Omega\chvec f\cdot\chvec v\,dx+
\int_\Gamma\chvec t\cdot\chvec v\,ds
\qquad\forall\chvec v\in\widehat{\chvec H}^1(\Omega).
\]
但~\eqref{eq:ch11-rigid-motion-strain-zero} 与相容性条件蕴含
\[
\MathBookFormulaID{formula-ch11-pure-traction-rigid-motion-equality}
a(\chvec u^*,\chvec v)=\int_\Omega\chvec f\cdot\chvec v\,dx+
\int_\Gamma\chvec t\cdot\chvec v\,ds
\qquad\forall\chvec v\in\chvec{\mathrm{RM}}.
\]
由于 $\chvec H^1(\Omega)=\widehat{\chvec H}^1(\Omega)\oplus\chvec{\mathrm{RM}}$, $\chvec u^*$ 是~\eqref{eq:ch11-pure-traction-variational-problem} 的一个解.
\end{Proof}

在关于 $\Omega$ 边界的某些条件下, $\chvec f\in\chvec L^2(\Omega)$, $\chvec g=\chvec u_1|_{\Gamma_1}$, 其中 $\chvec u_1\in\chvec H^2(\Omega)$, 且 $\chvec t=\chvec u_2|_{\Gamma_2}$, 其中 $\chvec u_2\in\chvec H^1(\Omega)$, 这一事实蕴含~\eqref{eq:ch11-elasticity-weak-problem} 的解属于 $\chvec H^2(\Omega)$, 然后命题~\MBUnresolvedReference{prop:ch05-source-5-1-9}{5.1.9} 的技术表明, $\chvec u$ 实际满足边值问题~\eqref{eq:ch11-static-elasticity-equation}--\eqref{eq:ch11-traction-boundary-condition}. 例如, 当 $\partial\Omega$ 光滑且 $\overline{\Gamma}_1\cap\overline{\Gamma}_2=\varnothing$(参见 Valent 1988), 或 $\Omega$ 为凸多边形且 $\Gamma_1$ 或 $\Gamma_2$ 为空(参见 Grisvard 1986 和 1989)时, 情形如此. 此外, 在这些情形下(Vogelius 1983, Brenner 与 Sung 1992), 存在与 $(\mu,\lambda)\in[\mu_1,\mu_2]\times(0,\infty)$ 无关的正常数 $C$, 使得
\MBSourceItem{33}
\begin{equation}
\MathBookFormulaID{formula-ch11-elasticity-regularity-estimate}
\label{eq:ch11-elasticity-regularity-estimate}
\|\chvec u\|_{\chvec H^2(\Omega)}+
\lambda\|\operatorname{div}\chvec u\|_{H^1(\Omega)}
\leq C\left(\|\chvec f\|_{\chvec L^2(\Omega)}+
\|\chvec u_1\|_{\chvec H^2(\Omega)}+
\|\chvec u_2\|_{\chvec H^1(\Omega)}\right).
\end{equation}
